\section{Introduction}
\label{sec:intro}

Image steganography traditionally hides information by modifying a carrier, whereas coverless approaches communicate through selection, indexing, semantic mapping, or generation. Selection preserves the original pixels, shifting the principal security question toward the distribution of transmitted images and the reliability of their labels after channel perturbations. Recent surveys accordingly separate cover-selection design from detectability and emphasize evaluation criteria that reflect the selected-carrier distribution \cite{azam2025coversurvey,kombrink2025survey}.

A finite shared index adds a second protocol constraint. Every protected message induces a concrete multiplicity demand over codebook symbols. Robust communication therefore requires enough qualified carrier groups for every requested symbol across the complete protected envelope. Stability-aware quantization and representative selection improve robust cover choice \cite{guo2026capacity}, while restoration and deep-hashing methods strengthen geometric robustness \cite{meng2025restoration,tan2026hashing}. NARCIS complements these directions by placing finite-index feasibility, authenticated message recovery, and selection detectability inside one end-to-end evaluation unit.

NARCIS couples a channel-aware quantile codebook with complementary five-cover groups whose majority-failure signatures are characterized on calibration transformations. A proportional-deficit scheduler preserves the natural signature mixture of each label bank, while an exact cyclic keyed Gray mapping balances coded symbols across visual clusters over every full $K$-session cycle. AES-GCM binds payload and session state, Reed--Solomon coding absorbs residual byte errors, and authenticated plaintext recovery becomes the final success criterion. The scientific contribution is this integrated finite-index construction and its externally frozen evaluation rather than any individual standard primitive.

The validation separates development evidence from external evidence. BOSSBase 1.01 \cite{bas2011boss} supports design and operating-point selection. The final protocol then evaluates Caltech-101 \cite{feifei2004caltech101,caltech101data} with $K=8$, five covers per coded symbol, 128 Reed--Solomon parity bytes, and 30 authenticated sessions per partition. The external campaign combines twelve calibration transformations, eight unseen holdout transformations, complete-index usage measurements, and repeated selection-steganalysis with SRM-lite, GLCM, and CNN detectors. A frozen execution of DiffStega \cite{yang2024diffstega} supplies an executable recent reference under its native reconstruction metrics.

The resulting evidence is broad across robustness, usage, and detectability. The BOSSBase development lineage achieves 1,800/1,800 calibration recoveries over five seeds, with 7,995--8,000 unique covers exercised per seed and mean repeated-session detector AUCs close to chance. On the external Caltech partitions, calibration reaches 1,800/1,800 authenticated recoveries and exercises all 7,000 covers in every partition. Unseen holdouts reach 1,163/1,200 authenticated recoveries (96.92\%), with the strongest 12\% crop defining the observed channel boundary. Final mean macro-AUCs are 0.504 for SRM-lite, 0.513 for GLCM, and 0.527 for the CNN. DiffStega completes all 100 official UniStega cases, with independently reproduced correct-recovery PSNR of 23.274 dB compared with 23.290 dB in the reference paper. Together, these results support strong bounded channel robustness and low measurable selection leakage under the evaluated threat model.

\section{Related Work and Positioning}
\label{sec:related}

\subsection{Selection, retrieval, and robust coverless communication}

Natural-image coverless schemes use block statistics, retrieval descriptors, object semantics, hashing, and learned features to map messages to carriers. Multi-object recognition \cite{luo2021multiobject} and joint selection/generation designs such as JoCS \cite{ren2024robust} illustrate the diversity of communication objects and mapping rules. Guo and Ping \cite{guo2026capacity} are especially close in spirit: their stability-aware pseudo-Zernike quantization and representative selection explicitly balance robustness and nominal capacity. NARCIS focuses on the complementary finite-index question: whether the complete authenticated coded demand can be realized by a fixed shared index and recovered at the plaintext level.

Geometric perturbations remain a central robustness challenge. Meng et al. \cite{meng2025restoration} restore attacked images before the underlying coverless decoder, while Tan et al. \cite{tan2026hashing} use deep unsupervised hashing to improve resistance to geometric attacks. NARCIS incorporates calibration transformations directly into codebook and group-bank construction, while eight stronger/intermediate values remain holdouts. The observed 12\% crop boundary therefore becomes an explicit external-validity result.

\subsection{Generative steganography and multimedia security}

Generative approaches communicate through synthesized content. DiffStega reconstructs a secret image through prompt- and password-conditioned diffusion \cite{yang2024diffstega}. Guidance-feature-distribution steganography embeds information in a generative feature distribution \cite{sun2024guidance}, and progressive generative steganography in TOMM jointly studies capacity, extraction, and detectability during high-resolution generation \cite{zhou2025pgs}. NARCIS operates on an indexed natural-image carrier set and measures exact authenticated byte recovery. The article therefore reports each family with its native information object and metrics.

Authentication also has several multimedia meanings. FaceSigns \cite{neekhara2024facesigns} embeds semi-fragile authentication information into media. NARCIS keeps indexed covers unchanged and authenticates the protected payload and session state, providing a complementary carrier-preserving authentication model.

\subsection{Steganalysis of selection policies}

Selection steganalysis tests whether the emitted subset differs statistically from the available carrier population. CNN-based image steganalysis can exploit subtle distributional differences \cite{li2023siairnet}; adversarial embedding work likewise shows that the evaluated steganalyst materially shapes the security conclusion \cite{wang2024adversarial}. Cao, Wang, and Zhang analyze cover-selection steganography directly and show that texture statistics can expose selection bias \cite{cao2025glcm}. NARCIS therefore evaluates the selected-vs.-remainder distinction with residual, texture, and learned detectors under image-disjoint train/test splits.

\begin{table}[t]
\caption{Positioning of representative recent approaches. Each row retains the semantics of its native communication object and evaluation protocol.}
\label{tab:positioning}
\small
\begin{tabularx}{\textwidth}{@{}lXXX@{}}
\toprule
Approach & Carrier / information object & Robustness mechanism & Security/evaluation emphasis \\
\midrule
Guo--Ping \cite{guo2026capacity} & Selected natural images / symbols & Stable PZM quantization and representative selection & Robustness and nominal capacity \\
Meng et al. \cite{meng2025restoration} & Existing coverless schemes & Attack classification and image restoration & Geometric robustness \\
DiffStega \cite{yang2024diffstega} & Generated images / secret image & Diffusion inversion/reconstruction & Perceptual recovery, password sensitivity, steganalysis \\
NARCIS & Unchanged indexed natural images / authenticated byte string & Calibration-locked codebook, complementary group bank, RS128 & Finite-index feasibility, exact authenticated recovery, selection detectability \\
\bottomrule
\end{tabularx}
\end{table}
